\subsection{Durchlauf der Tests}\label{subsec:durchlauf_der_tests}
Nachdem die Testdaten erstellt wurden, wurden die Tests mit verschiedenen Konfigurationen des Algorithmus durchgeführt. Dabei wurden verschiedene Werte für die Schwellenwerte der Prominenz und Dominanz getestet, um zu sehen, wie sich diese auf die Genauigkeit der Gipfelerkennung auswirken. Dabei ist aufgefallen, dass die Datenqualität aprupt abnimmt, wenn der Schwellenwert für die Prominenz und der Dominanz zu neidrig angesetzt werden. Um einerseits den Gipfelfinder in schwerem Gebiet zu Testen und gleichzeitig die Datenqualität zu erhalten, wurde bei keinem Test der Schwellenwert für die Prominenz unter 100m und der Schwellenwert für die Dominanz unter 1000m gesetzt. Es wurden aber auch Tests mit höheren Schwellenwerten durchgeführt, um zu sehen, wie sich diese auf die Genauigkeit auswirken.
